\documentclass[a4paper,12pt]{article}

\usepackage[utf8]{inputenc}
\usepackage{amsmath}
\usepackage{amssymb}
\usepackage{amsthm}
\usepackage{mathtools}
\usepackage[strings]{underscore}
\usepackage[margin=1in]{geometry}
\usepackage{hyperref}
\usepackage[numbers]{natbib}
\usepackage{xcolor}
\usepackage{hyperref}
\hypersetup{
	colorlinks,
	linkcolor={blue!100!black},
	citecolor={blue!100!black},
	urlcolor={blue!100!black}
}
\title{Foundations of General Systems Theory\\
\small \scshape A Survey of Formal Definitions}
\author{\scshape Igor Strozzi \\ \small \itshape Applied Mathematics Department \\ \small Federal University of Rio de Janeiro}
\date{}

\begin{document}

\maketitle

\begin{abstract}
This document presents a systematic compilation of foundational definitions of systems and major theoretical insights drawn from the mathematical literature on General Systems Theory. The conceptual foundations are traced through the seminal contributions of Mesarovic, Takahara, Yi Lin, and other key figures in the formalization of systems theory within set-theoretic and logical frameworks.
\end{abstract}

\tableofcontents

\newpage

\section{General Systems Theory}

\subsection{Foundational Definitions of a System}

The mathematical formalization of systems theory emerges from diverse epistemological traditions, each seeking to articulate the fundamental nature of systemic relationships through rigorous set-theoretic frameworks. We present here the principal definitions that constitute the conceptual architecture of General Systems Theory.

\subsubsection{The Verbal-Conceptual Definition}

Mesarovic \citep{Mesarovic1972} articulates the primordial intuition underlying all systems thinking:

\begin{quote}
\emph{A system is (or reflects the existence of) a relationship between the objects (of study).}
\end{quote}

This formulation reveals the essential phenomenological stance: systems are not \emph{things-in-themselves}, but rather patterns of interrelation discerned through observational praxis. The definition achieves generality precisely through its abstraction from particular physical instantiations, concerning itself instead with the structural configurations of relationships.

Mesarovic further refines this characterization to emphasize its teleological dimension:

\begin{quote}
\emph{A system is a relationship among objects described (or, specified, defined) in terms of information-processing and decision-making concepts.}
\end{quote}

\subsubsection{The Set-Theoretic Formalization}

The transition from verbal description to mathematical precision constitutes the foundational gesture of modern systems theory. Mesarovic \citep{Mesarovic1972} proposes the canonical set-theoretic definition:

\begin{align*}
S \subseteq V_1 \times \cdots \times V_n
\end{align*}

where $\times$ denotes the Cartesian product. The coordinate sets $V_1, \ldots, V_n$ are designated as \emph{system objects}. This formulation encapsulates several profound insights:

\begin{itemize}
\item An object $V_i$ represents the totality of all possible appearances of an attribute of the phenomenon under consideration.
\item The system $S$ specifies which combinations of object appearances can occur simultaneously when objects are interrelated to form the given system.
\item The relation $S$ is not merely an arbitrary subset, but rather embodies the structural constraints that govern the co-variation of system attributes.
\end{itemize}

Mesarovic and Takahara further develop this framework, emphasizing that systems are defined on the set-theoretic level with minimal mathematical structure \citep{Mesarovic1975}. They write:

\begin{quote}
A family of sets is given, $\mathcal{V} = \{V_i : i \in I\}$ where $I$ is the index set, and a system, defined on $\mathcal{V}$, is a proper subset of $\times \mathcal{V}$:
\begin{align*}
S \subseteq \times \{V_i : i \in I\}
\end{align*}
\end{quote}

\subsubsection{The Input-Output System}

For systems exhibiting causal structure, the formalism specializes to an input-output configuration. The most widely studied case considers a system consisting of two objects \citep{Mesarovic1975}:

\begin{align*}
S \subseteq X \times Y
\end{align*}

where $X$ represents the input object (stimuli) and $Y$ represents the output object (responses). This partitioning introduces the fundamental distinction between cause and effect, transforming the abstract relational structure into a framework amenable to dynamical analysis.

When $S$ is a function, i.e., with every input $x \in X$ there is associated a unique output, the system is designated a \emph{functional system}:

\begin{align*}
S: X \to Y
\end{align*}

This functional characterization permits the introduction of \emph{internal mechanisms}---constructive specifications that determine output generation for any given input through differential equations, transformation tables, or algorithmic procedures.

\subsubsection{The Multirelation Approach}

Yi Lin \citep{Lin1999} introduces an alternative formalization emphasizing the multiplicity of relations within a system. A system is defined as an ordered pair:

\begin{align*}
S = (M, R)
\end{align*}

where $M$ is a set of \emph{objects} and $R$ is a set of \emph{relations} on $M$. For each relation $r \in R$, there exists an ordinal number $n = n(r)$ such that $r \subseteq M^n$, where $M^n$ denotes the $n$-fold Cartesian product of $M$.

This formulation explicitly acknowledges that systems typically exhibit multiple simultaneous relationships of varying arities, capturing the richness of interconnection patterns that characterize complex phenomena.

\subsubsection{Bunge's Compositional Definition}

Bunge introduces environmental considerations into the system definition \citep[as cited in][]{Lin1999,Bunge1979}. Given a nonempty set $T$, an ordered triple $W = (C, E, S)$ constitutes a system over $T$ if and only if:

\begin{itemize}
\item $C$ and $E$ are mutually disjoint subsets of $T$ (composition and environment, respectively)
\item $S$ is a nonempty set of relations on $C \cup E$
\end{itemize}

This tripartite structure distinguishes internal composition from external environment, while acknowledging that systemic relations may span both domains.
\subsection{Major Theoretical Insights}

\subsubsection{The Formalization Methodology}

A central methodological principle articulated by Mesarovic and Takahara involves the systematic progression from minimal to enriched mathematical structure \citep{Mesarovic1975}:

\begin{enumerate}
\item \textbf{Formalization}: Basic systems concepts are introduced through precise mathematical definitions using minimal structure---as few axioms as correct interpretation permits.
\item \textbf{Structural Enrichment}: Mathematical theory develops by adding structure as needed for investigating various systems properties, thereby revealing which properties are fundamental and which assumptions are minimally necessary.
\end{enumerate}

This approach manifests deep epistemological commitments regarding the relationship between mathematical abstraction and phenomenological adequacy.

\subsubsection{Terminal Systems versus Goal-Seeking Systems}

Mesarovic \citep{Mesarovic1972} introduces a fundamental bifurcation in systems description:

\begin{itemize}
\item \textbf{Terminal Systems (Causal Systems)}: Systems considered as processing information flows, with detailed specification given in terms of internal mechanisms. The focus is on \emph{how} inputs transform into outputs.

\item \textbf{Goal-Seeking Systems}: Systems pursuing objectives or purposes, with outputs described as goal-seeking responses to inputs. These encompass control systems, decision-making systems, and adaptive systems.
\end{itemize}

For goal-seeking systems, the state-space representation becomes essential. If $X$ and $Y$ are time functions and $Z$ is a set of state functions, the system $S \subseteq X \times Y$ is made functional through introduction of state:

\begin{align*}
S: Z \times X \to Y
\end{align*}

State evolution is governed by transition functions, commonly expressed as differential or difference equations.

\subsubsection{Hierarchical and Multi-Level Systems}

The analysis of large-scale complex systems necessitates hierarchical decomposition \citep{Mesarovic1972}. Mesarovic articulates several principles:

\begin{itemize}
\item \textbf{Stratified Description}: Complex systems naturally decompose into subsystems across multiple levels. At any particular level, subsystems are considered in isolation; at higher levels, they are incorporated into larger subsystems.

\item \textbf{Multi-Level Multi-Goal Systems}: The most sophisticated organizational structures contain interacting subsystems with distinct (potentially conflicting) goals, arranged in hierarchical supremacy relations where higher-level subsystems influence, condition, or control goal-seeking activities of lower-level subsystems.
\end{itemize}

The concept of hierarchy emerges not as an arbitrary organizational principle but as an intrinsic feature of large-scale systems, intimately related to the very notion of complexity.

\subsubsection{System Models and Structures}

Takahashi and Takahara \citep{Takahashi1995} emphasize the critical distinction between \emph{system models} and \emph{system structures}:

\begin{itemize}
\item A \textbf{system model} is a mathematical structure $\mathcal{M}$ serving as the formal representation of a system.

\item The \textbf{structure of a system} is defined as an ordered pair $(\mathcal{L}(\mathcal{M}), \Sigma)$, where:
\begin{itemize}
\item $\mathcal{L}(\mathcal{M})$ is the formal language for describing the system model
\item $\Sigma$ is a set of axioms (sentences in $\mathcal{L}(\mathcal{M})$) specifying how elements interact
\end{itemize}
\end{itemize}

This distinction reveals that a single system model may possess multiple structures, corresponding to different aspects or levels of analysis. The structure captures the modeler's recognition of which interactions are salient, rather than an objective property of the system itself.

\subsubsection{Universality and Expressive Power}

A remarkable insight concerns the expressive adequacy of the set-theoretic formulation. Mesarovic demonstrates that even verbal descriptions can be formalized as systems \citep{Mesarovic1972}:

\begin{quote}
Every verbal statement possesses two basic linguistic categories: \emph{nouns} (terms) denoting objects of concern, and \emph{functors} denoting relations between terms. If terms range over sets satisfying axiomatic set theory, functors represent relations on these sets. To every collection of statements representing a theory of phenomena, there can be assigned a mathematical model---a system in the sense of the set-theoretic definition.
\end{quote}

The limitations imposed by General Systems Theory derive solely from requiring that the range of any undetermined term satisfies axiomatic set theory, which essentially amounts to avoiding logical paradoxes. The inadequacy of any particular system model reflects limitations in our phenomenological theory rather than constraints inherent to the formalism.

\subsubsection{Fundamental Objects and Layer Structures}

Yi Lin \citep{Lin1999} establishes that every system $S$ possesses a unique set $M(S)$ consisting of all \emph{fundamental objects}---objects which are not themselves systems. This leads to a layered conception:

\begin{itemize}
\item Systems exhibit hierarchical object structures, where higher-level objects may themselves be systems
\item The Axiom of Regularity ensures that chains of object-subsystem relations must be finite
\item Every system can be decomposed into fundamental objects through iterative extraction of non-systemic elements
\end{itemize}

This structural characterization provides rigorous foundations for analyzing part-whole relations and systemic decomposition.

\subsubsection{The Primacy of Relationships over Mechanisms}

A philosophical cornerstone articulated across multiple sources emphasizes that systems are defined through \emph{observed features and their relationships} rather than underlying mechanisms \citep{Mesarovic1975}:

\begin{quote}
This accords with the nature of the systems field and its concern with organization and interrelationships of components into an overall system, rather than with specific mechanisms within a given phenomenological framework.
\end{quote}

This principle establishes systems theory as fundamentally concerned with structural patterns that persist across diverse physical, biological, social, and abstract domains---a genuinely transdisciplinary framework for analyzing organized complexity.

\bibliographystyle{plain}
\bibliography{references/bibliography.bib}

\end{document}
